% ============================================================
%  MAINTAINER'S GUIDE
%  Bug-Bite Classification via Tropical Cyclone Transfer Learning
%  Braude College of Engineering — Software Engineering Capstone
%
%  Usage in Overleaf:
%    \appendix
%    % ============================================================
%  MAINTAINER'S GUIDE
%  Bug-Bite Classification via Tropical Cyclone Transfer Learning
%  Braude College of Engineering — Software Engineering Capstone
%
%  Usage in Overleaf:
%    \appendix
%    % ============================================================
%  MAINTAINER'S GUIDE
%  Bug-Bite Classification via Tropical Cyclone Transfer Learning
%  Braude College of Engineering — Software Engineering Capstone
%
%  Usage in Overleaf:
%    \appendix
%    % ============================================================
%  MAINTAINER'S GUIDE
%  Bug-Bite Classification via Tropical Cyclone Transfer Learning
%  Braude College of Engineering — Software Engineering Capstone
%
%  Usage in Overleaf:
%    \appendix
%    \input{maintainer_guide}
%
%  Or compile standalone:
%    Uncomment the preamble block below, and add \end{document} at the end.
% ============================================================

% --- Standalone preamble (remove if including in main document) ----------
\documentclass[12pt,a4paper]{article}
\usepackage[utf8]{inputenc}
\usepackage[T1]{fontenc}
\usepackage{geometry}
\geometry{margin=2.5cm}
\usepackage{hyperref}
\usepackage{listings}
\usepackage{xcolor}
\usepackage{booktabs}
\usepackage{array}
\usepackage{enumitem}
\usepackage{parskip}
\lstset{basicstyle=\ttfamily\small, breaklines=true, frame=single,
        backgroundcolor=\color{gray!8}}
\begin{document}
\title{Maintainer's Guide\\
       \large Bug-Bite Classification via Tropical Cyclone Transfer Learning\\
       Braude College of Engineering --- Software Engineering Capstone}
\author{Ishay Yulzary, Nathan Trostianitser}
\date{}
\maketitle
\tableofcontents
\newpage
% -------------------------------------------------------------------------

\section{Maintainer's Guide}
\label{app:maintainer-guide}

This guide is addressed to a software engineer or researcher responsible for
keeping this system operational, extending the research, or advancing the
prototype toward a deployable product. It covers the operating environment,
installation of all custom software, the training pipeline, and procedures
for updating and extending the system.

% ============================================================
\subsection{System Overview}
\label{sec:mg-overview}
% ============================================================

\subsubsection{What the System Does}

The system classifies a photograph of a bug bite into one of five categories:
ant bite, bed bug bite, mosquito bite, spider bite, or tick/flea bite. It does
this by running the photograph through an ensemble of three neural network
models and produces visual explanations showing which image regions most
influenced the prediction.

The inference interface is \texttt{Stacked\_Model\_Colab.ipynb}, a
self-contained Jupyter notebook designed for Google Colab. The full training
and evaluation pipeline runs via \texttt{miscellaneous\_code/run\_experiments.py}
on a dedicated GPU machine.


% ============================================================
\subsection{Repository Structure}
\label{sec:mg-repo}
% ============================================================

\begin{description}
    \item[\texttt{notebooks/Multiclass\_Classification.ipynb}]
          Primary experimental notebook. Covers all four pipeline stages:
          control training, TC pre-training, feature map and GradCAM
          evaluation, and the full XAI suite. Intended for interactive
          exploration or one-shot reproduction on a GPU instance.

    \item[\texttt{Stacked\_Model\_Colab.ipynb}]
          Self-contained inference and visualisation notebook for Google Colab
          free tier. Loads trained weights from Google Drive, runs ensemble
          inference on a single image, and produces side-by-side Control vs.\
          TC GradCAM and LIME outputs.

    \item[\texttt{scripts/}]
          Standalone Python scripts, one per pipeline stage, designed for
          unattended server execution via \texttt{run\_experiments.py}.
          \begin{itemize}
              \item \texttt{shared.py} --- shared imports, dataset paths,
                    metrics helpers, and GradCAM utilities.
              \item \texttt{01\_train\_control.py} --- trains the Control
                    (ImageNet $\rightarrow$ bug-bite) models.
              \item \texttt{02\_train\_tc.py} --- runs cyclone pre-training
                    then TC $\rightarrow$ bug-bite fine-tuning.
              \item \texttt{03\_evaluate.py} --- generates feature map
                    comparisons and GradCAM overlays.
              \item \texttt{04\_xai.py} --- computes ensemble metrics,
                    silhouette scores, t-SNE projections, LIME, SHAP, and
                    DiCE outputs.
          \end{itemize}

    \item[\texttt{miscellaneous\_code/pytorch\_utils.py}]
          Custom PyTorch training library imported by all scripts and the
          main notebook. See Section~\ref{sec:mg-pytorch-utils}.

    \item[\texttt{miscellaneous\_code/run\_experiments.py}]
          Automated hyperparameter sweep runner.
          See Section~\ref{sec:mg-sweep}.

    \item[\texttt{miscellaneous\_code/fetch\_cyclone\_dataset.py}]
          Downloads and labels the HURSAT-B1 cyclone dataset from NOAA.
          See Section~\ref{sec:mg-data-cyclone}.

    \item[\texttt{miscellaneous\_code/cyclone\_preprocessing.py}]
          Removes grid-line artefacts introduced by HURSAT-B1 image tiling
          before training.

    \item[\texttt{Bug-Data/Multiclass\_Bug\_data/}]
          Bug-bite dataset: \texttt{train/\{class\}/} and
          \texttt{val/\{class\}/} subdirectories, one per class.

    \item[\texttt{results/}]
          Auto-generated by the sweep runner. Contains checkpoints,
          per-run metric JSONs, XAI outputs, and plots.

    \item[\texttt{deploy/}]
          Legacy Streamlit demo (binary classifier, outdated). Not part
          of the active pipeline; retained for reference only.
\end{description}

% ============================================================
\subsection{Operating Environment}
\label{sec:mg-env}
% ============================================================

\subsubsection{Hardware Requirements}

\textbf{A dedicated GPU is required for training.} The full 63-configuration
sweep was run on the following hardware:

\begin{itemize}
    \item \textbf{CPU:} Intel Core i7-9700KF @ 5.10\,GHz all-core OC
          (8 cores / 8 threads)
    \item \textbf{RAM:} 32\,GB
    \item \textbf{GPU:} NVIDIA RTX A5000 (24\,GB GDDR6 VRAM)
\end{itemize}

The full sweep took approximately 100--130 hours of continuous GPU time on
this system. A GPU with at least 16\,GB VRAM is recommended; the default
batch size of 4 allows training on GPUs with as little as 8\,GB VRAM at the
cost of slower convergence. InceptionV3 phase-1 training may require a
reduced batch size on GPUs below 24\,GB VRAM (see
Section~\ref{sec:mg-troubleshoot-oom}).

Google Colab free tier is not suitable for the full sweep due to session time
limits. A Colab Pro instance with a high-RAM GPU runtime is the minimum viable
option for partial retraining. Inference and XAI visualisation can be run on
CPU or a smaller GPU.

\subsubsection{Software Dependencies}

\textbf{Inference (Colab):} The inference notebook installs its own
dependencies at runtime in Section 1. No manual installation is required.

\textbf{Training (local or server):}

\begin{lstlisting}[language=bash]
# Core training dependencies
pip install torch torchvision timm numpy Pillow opencv-python \
            scikit-learn matplotlib tqdm

# Optional XAI packages
pip install shap lime scikit-image dice-ml

# Cyclone dataset fetch
pip install netCDF4 requests pandas tqdm
\end{lstlisting}

There is no pinned \texttt{requirements.txt} for the training environment.
The project was developed and tested with PyTorch 2.x and timm 0.9.x.
Contact \texttt{ishay.yulzary@gmail.com} for the exact environment
specification used during the published experiments.

\subsubsection{Dataset Paths}

Dataset paths are configured at the top of
\texttt{miscellaneous\_code/run\_experiments.py} and in
\texttt{scripts/shared.py}. They default to WSL native paths:

\begin{lstlisting}[language=Python]
BUG_TRAIN = '/home/test/bug_data/train'
BUG_VAL   = '/home/test/bug_data/val'
TC_TRAIN  = '/home/test/cyclone_data_split/train'
TC_VAL    = '/home/test/cyclone_data_split/val'
\end{lstlisting}

Update these constants to match your environment before running training.

% ============================================================
\subsection{Model Weights and Artifacts}
\label{sec:mg-weights}
% ============================================================

Trained model weights are not stored in the repository. Two archives are
available from \texttt{ishay.yulzary@gmail.com}:

\begin{itemize}
    \item \texttt{project\_book\_results.tar} --- the selected best-performing
          checkpoints used in the inference demo (one Control config, one TC
          config, three architectures each).
    \item \textbf{Complete sweep archive} --- all 70 experiment configurations
          (63 TC + 7 Control seeds), covering every trained checkpoint from
          the published experiments. Request this for full evaluation
          reproduction or sweep extension.
\end{itemize}

Weight subdirectories within the demo archive:
\begin{itemize}
    \item Control: \texttt{project\_book\_results/s3\_ls0.20\_lr1.0e-05/}
    \item TC: \texttt{project\_book\_results/s6\_ls0.10\_lr5.0e-06/}
\end{itemize}

\textbf{Updating the weights archive after retraining:}
New weight files are written to \texttt{results/checkpoints/}. Package the
relevant checkpoint directories into a new
\\\texttt{project\_book\_results.tar} and upload it to Google Drive, replacing
the previous file.


% ============================================================
\subsection{Troubleshooting}
\label{sec:mg-troubleshoot}
% ============================================================

\subsubsection{GPU out-of-memory error during training}
\label{sec:mg-troubleshoot-oom}

\textbf{Symptom:} \texttt{torch.cuda.OutOfMemoryError} during
\texttt{run\_experiments.py}.

\textbf{Resolution:}
\begin{enumerate}
    \item Reduce the batch size for the affected architecture in the
          \texttt{MODELS} list in \texttt{run\_experiments.py}.
          Relevant columns: \texttt{cyc\_p1\_bs} (cyclone phase-1 batch
          size) and \texttt{bug\_bs} (bug-bite batch size).
    \item For InceptionV3, the phase-1 batch size may need to be reduced
          to 64 or below on GPUs with less than 24\,GB VRAM.
    \item Re-run the sweep. Configurations that completed before the
          crash are not repeated.
\end{enumerate}

\subsubsection{NOAA download failures during cyclone dataset fetch}

\textbf{Symptom:} Repeated warnings such as
\begin{lstlisting}
! Index fetch failed for <year>: ...
\end{lstlisting}

\textbf{Resolution:}
\begin{enumerate}
    \item Re-run the fetch script. Downloads are resumable; already-saved
          images are not re-downloaded.
    \item Reduce the worker count if the server rate-limits the connection:
\begin{lstlisting}[language=bash]
python miscellaneous_code/fetch_cyclone_dataset.py --workers 2
\end{lstlisting}
    \item If NOAA servers are unavailable, contact
          \texttt{ishay.yulzary@gmail.com} to obtain the pre-built dataset
          archive directly.
\end{enumerate}

\subsubsection{Training sweep shows no progress after resuming}

\textbf{Symptom:} Re-running \texttt{run\_experiments.py} prints
\texttt{All configs already completed} but results appear incomplete.

\textbf{Resolution:}
\begin{enumerate}
    \item Run \texttt{python miscellaneous\_code/run\_experiments.py
          --status} to view the current state of each configuration.
    \item If all 63 configurations report \texttt{completed} but results
          are missing, \\\texttt{results/experiment\_results.json} may be
          corrupted. Delete it and re-run; the sweep will re-evaluate all
          existing checkpoints without retraining.
    \item If fewer than 63 configurations are completed, pending ones
          resume automatically on the next run.
\end{enumerate}

% ============================================================
\subsection{Contacts and Resources}
\label{sec:mg-contacts}
% ============================================================

\begin{description}
    \item[Primary contact] Ishay Yulzary ---
          \texttt{ishay.yulzary@gmail.com}.
          Contact for: weight archive access, dataset access, environment
          specification, and questions not covered in this guide.

    \item[Repository] Source code, notebooks, and bug-bite dataset are
          maintained in the project GitHub repository.

    \item[Data sources] The cyclone dataset is derived from:
          \begin{itemize}
              \item HURSAT-B1:
                    \\\url{https://www.ncei.noaa.gov/products/hurricane-satellite-data}
              \item IBTrACS:
                    \\\url{https://www.ncei.noaa.gov/products/international-best-track-archive}
          \end{itemize}
          Both are freely available from NOAA.

    \item[Model library] Architectures are loaded via the \texttt{timm}
          library (\url{https://github.com/huggingface/pytorch-image-models}).
          Pre-trained ImageNet weights are downloaded automatically by
          \texttt{timm} on first use.
\end{description}

% --- End of appendix (remove if standalone) ---
\end{document}

%
%  Or compile standalone:
%    Uncomment the preamble block below, and add \end{document} at the end.
% ============================================================

% --- Standalone preamble (remove if including in main document) ----------
\documentclass[12pt,a4paper]{article}
\usepackage[utf8]{inputenc}
\usepackage[T1]{fontenc}
\usepackage{geometry}
\geometry{margin=2.5cm}
\usepackage{hyperref}
\usepackage{listings}
\usepackage{xcolor}
\usepackage{booktabs}
\usepackage{array}
\usepackage{enumitem}
\usepackage{parskip}
\lstset{basicstyle=\ttfamily\small, breaklines=true, frame=single,
        backgroundcolor=\color{gray!8}}
\begin{document}
\title{Maintainer's Guide\\
       \large Bug-Bite Classification via Tropical Cyclone Transfer Learning\\
       Braude College of Engineering --- Software Engineering Capstone}
\author{Ishay Yulzary, Nathan Trostianitser}
\date{}
\maketitle
\tableofcontents
\newpage
% -------------------------------------------------------------------------

\section{Maintainer's Guide}
\label{app:maintainer-guide}

This guide is addressed to a software engineer or researcher responsible for
keeping this system operational, extending the research, or advancing the
prototype toward a deployable product. It covers the operating environment,
installation of all custom software, the training pipeline, and procedures
for updating and extending the system.

% ============================================================
\subsection{System Overview}
\label{sec:mg-overview}
% ============================================================

\subsubsection{What the System Does}

The system classifies a photograph of a bug bite into one of five categories:
ant bite, bed bug bite, mosquito bite, spider bite, or tick/flea bite. It does
this by running the photograph through an ensemble of three neural network
models and produces visual explanations showing which image regions most
influenced the prediction.

The inference interface is \texttt{Stacked\_Model\_Colab.ipynb}, a
self-contained Jupyter notebook designed for Google Colab. The full training
and evaluation pipeline runs via \texttt{miscellaneous\_code/run\_experiments.py}
on a dedicated GPU machine.


% ============================================================
\subsection{Repository Structure}
\label{sec:mg-repo}
% ============================================================

\begin{description}
    \item[\texttt{notebooks/Multiclass\_Classification.ipynb}]
          Primary experimental notebook. Covers all four pipeline stages:
          control training, TC pre-training, feature map and GradCAM
          evaluation, and the full XAI suite. Intended for interactive
          exploration or one-shot reproduction on a GPU instance.

    \item[\texttt{Stacked\_Model\_Colab.ipynb}]
          Self-contained inference and visualisation notebook for Google Colab
          free tier. Loads trained weights from Google Drive, runs ensemble
          inference on a single image, and produces side-by-side Control vs.\
          TC GradCAM and LIME outputs.

    \item[\texttt{scripts/}]
          Standalone Python scripts, one per pipeline stage, designed for
          unattended server execution via \texttt{run\_experiments.py}.
          \begin{itemize}
              \item \texttt{shared.py} --- shared imports, dataset paths,
                    metrics helpers, and GradCAM utilities.
              \item \texttt{01\_train\_control.py} --- trains the Control
                    (ImageNet $\rightarrow$ bug-bite) models.
              \item \texttt{02\_train\_tc.py} --- runs cyclone pre-training
                    then TC $\rightarrow$ bug-bite fine-tuning.
              \item \texttt{03\_evaluate.py} --- generates feature map
                    comparisons and GradCAM overlays.
              \item \texttt{04\_xai.py} --- computes ensemble metrics,
                    silhouette scores, t-SNE projections, LIME, SHAP, and
                    DiCE outputs.
          \end{itemize}

    \item[\texttt{miscellaneous\_code/pytorch\_utils.py}]
          Custom PyTorch training library imported by all scripts and the
          main notebook. See Section~\ref{sec:mg-pytorch-utils}.

    \item[\texttt{miscellaneous\_code/run\_experiments.py}]
          Automated hyperparameter sweep runner.
          See Section~\ref{sec:mg-sweep}.

    \item[\texttt{miscellaneous\_code/fetch\_cyclone\_dataset.py}]
          Downloads and labels the HURSAT-B1 cyclone dataset from NOAA.
          See Section~\ref{sec:mg-data-cyclone}.

    \item[\texttt{miscellaneous\_code/cyclone\_preprocessing.py}]
          Removes grid-line artefacts introduced by HURSAT-B1 image tiling
          before training.

    \item[\texttt{Bug-Data/Multiclass\_Bug\_data/}]
          Bug-bite dataset: \texttt{train/\{class\}/} and
          \texttt{val/\{class\}/} subdirectories, one per class.

    \item[\texttt{results/}]
          Auto-generated by the sweep runner. Contains checkpoints,
          per-run metric JSONs, XAI outputs, and plots.

    \item[\texttt{deploy/}]
          Legacy Streamlit demo (binary classifier, outdated). Not part
          of the active pipeline; retained for reference only.
\end{description}

% ============================================================
\subsection{Operating Environment}
\label{sec:mg-env}
% ============================================================

\subsubsection{Hardware Requirements}

\textbf{A dedicated GPU is required for training.} The full 63-configuration
sweep was run on the following hardware:

\begin{itemize}
    \item \textbf{CPU:} Intel Core i7-9700KF @ 5.10\,GHz all-core OC
          (8 cores / 8 threads)
    \item \textbf{RAM:} 32\,GB
    \item \textbf{GPU:} NVIDIA RTX A5000 (24\,GB GDDR6 VRAM)
\end{itemize}

The full sweep took approximately 100--130 hours of continuous GPU time on
this system. A GPU with at least 16\,GB VRAM is recommended; the default
batch size of 4 allows training on GPUs with as little as 8\,GB VRAM at the
cost of slower convergence. InceptionV3 phase-1 training may require a
reduced batch size on GPUs below 24\,GB VRAM (see
Section~\ref{sec:mg-troubleshoot-oom}).

Google Colab free tier is not suitable for the full sweep due to session time
limits. A Colab Pro instance with a high-RAM GPU runtime is the minimum viable
option for partial retraining. Inference and XAI visualisation can be run on
CPU or a smaller GPU.

\subsubsection{Software Dependencies}

\textbf{Inference (Colab):} The inference notebook installs its own
dependencies at runtime in Section 1. No manual installation is required.

\textbf{Training (local or server):}

\begin{lstlisting}[language=bash]
# Core training dependencies
pip install torch torchvision timm numpy Pillow opencv-python \
            scikit-learn matplotlib tqdm

# Optional XAI packages
pip install shap lime scikit-image dice-ml

# Cyclone dataset fetch
pip install netCDF4 requests pandas tqdm
\end{lstlisting}

There is no pinned \texttt{requirements.txt} for the training environment.
The project was developed and tested with PyTorch 2.x and timm 0.9.x.
Contact \texttt{ishay.yulzary@gmail.com} for the exact environment
specification used during the published experiments.

\subsubsection{Dataset Paths}

Dataset paths are configured at the top of
\texttt{miscellaneous\_code/run\_experiments.py} and in
\texttt{scripts/shared.py}. They default to WSL native paths:

\begin{lstlisting}[language=Python]
BUG_TRAIN = '/home/test/bug_data/train'
BUG_VAL   = '/home/test/bug_data/val'
TC_TRAIN  = '/home/test/cyclone_data_split/train'
TC_VAL    = '/home/test/cyclone_data_split/val'
\end{lstlisting}

Update these constants to match your environment before running training.

% ============================================================
\subsection{Model Weights and Artifacts}
\label{sec:mg-weights}
% ============================================================

Trained model weights are not stored in the repository. Two archives are
available from \texttt{ishay.yulzary@gmail.com}:

\begin{itemize}
    \item \texttt{project\_book\_results.tar} --- the selected best-performing
          checkpoints used in the inference demo (one Control config, one TC
          config, three architectures each).
    \item \textbf{Complete sweep archive} --- all 70 experiment configurations
          (63 TC + 7 Control seeds), covering every trained checkpoint from
          the published experiments. Request this for full evaluation
          reproduction or sweep extension.
\end{itemize}

Weight subdirectories within the demo archive:
\begin{itemize}
    \item Control: \texttt{project\_book\_results/s3\_ls0.20\_lr1.0e-05/}
    \item TC: \texttt{project\_book\_results/s6\_ls0.10\_lr5.0e-06/}
\end{itemize}

\textbf{Updating the weights archive after retraining:}
New weight files are written to \texttt{results/checkpoints/}. Package the
relevant checkpoint directories into a new
\\\texttt{project\_book\_results.tar} and upload it to Google Drive, replacing
the previous file.


% ============================================================
\subsection{Troubleshooting}
\label{sec:mg-troubleshoot}
% ============================================================

\subsubsection{GPU out-of-memory error during training}
\label{sec:mg-troubleshoot-oom}

\textbf{Symptom:} \texttt{torch.cuda.OutOfMemoryError} during
\texttt{run\_experiments.py}.

\textbf{Resolution:}
\begin{enumerate}
    \item Reduce the batch size for the affected architecture in the
          \texttt{MODELS} list in \texttt{run\_experiments.py}.
          Relevant columns: \texttt{cyc\_p1\_bs} (cyclone phase-1 batch
          size) and \texttt{bug\_bs} (bug-bite batch size).
    \item For InceptionV3, the phase-1 batch size may need to be reduced
          to 64 or below on GPUs with less than 24\,GB VRAM.
    \item Re-run the sweep. Configurations that completed before the
          crash are not repeated.
\end{enumerate}

\subsubsection{NOAA download failures during cyclone dataset fetch}

\textbf{Symptom:} Repeated warnings such as
\begin{lstlisting}
! Index fetch failed for <year>: ...
\end{lstlisting}

\textbf{Resolution:}
\begin{enumerate}
    \item Re-run the fetch script. Downloads are resumable; already-saved
          images are not re-downloaded.
    \item Reduce the worker count if the server rate-limits the connection:
\begin{lstlisting}[language=bash]
python miscellaneous_code/fetch_cyclone_dataset.py --workers 2
\end{lstlisting}
    \item If NOAA servers are unavailable, contact
          \texttt{ishay.yulzary@gmail.com} to obtain the pre-built dataset
          archive directly.
\end{enumerate}

\subsubsection{Training sweep shows no progress after resuming}

\textbf{Symptom:} Re-running \texttt{run\_experiments.py} prints
\texttt{All configs already completed} but results appear incomplete.

\textbf{Resolution:}
\begin{enumerate}
    \item Run \texttt{python miscellaneous\_code/run\_experiments.py
          --status} to view the current state of each configuration.
    \item If all 63 configurations report \texttt{completed} but results
          are missing, \\\texttt{results/experiment\_results.json} may be
          corrupted. Delete it and re-run; the sweep will re-evaluate all
          existing checkpoints without retraining.
    \item If fewer than 63 configurations are completed, pending ones
          resume automatically on the next run.
\end{enumerate}

% ============================================================
\subsection{Contacts and Resources}
\label{sec:mg-contacts}
% ============================================================

\begin{description}
    \item[Primary contact] Ishay Yulzary ---
          \texttt{ishay.yulzary@gmail.com}.
          Contact for: weight archive access, dataset access, environment
          specification, and questions not covered in this guide.

    \item[Repository] Source code, notebooks, and bug-bite dataset are
          maintained in the project GitHub repository.

    \item[Data sources] The cyclone dataset is derived from:
          \begin{itemize}
              \item HURSAT-B1:
                    \\\url{https://www.ncei.noaa.gov/products/hurricane-satellite-data}
              \item IBTrACS:
                    \\\url{https://www.ncei.noaa.gov/products/international-best-track-archive}
          \end{itemize}
          Both are freely available from NOAA.

    \item[Model library] Architectures are loaded via the \texttt{timm}
          library (\url{https://github.com/huggingface/pytorch-image-models}).
          Pre-trained ImageNet weights are downloaded automatically by
          \texttt{timm} on first use.
\end{description}

% --- End of appendix (remove if standalone) ---
\end{document}

%
%  Or compile standalone:
%    Uncomment the preamble block below, and add \end{document} at the end.
% ============================================================

% --- Standalone preamble (remove if including in main document) ----------
\documentclass[12pt,a4paper]{article}
\usepackage[utf8]{inputenc}
\usepackage[T1]{fontenc}
\usepackage{geometry}
\geometry{margin=2.5cm}
\usepackage{hyperref}
\usepackage{listings}
\usepackage{xcolor}
\usepackage{booktabs}
\usepackage{array}
\usepackage{enumitem}
\usepackage{parskip}
\lstset{basicstyle=\ttfamily\small, breaklines=true, frame=single,
        backgroundcolor=\color{gray!8}}
\begin{document}
\title{Maintainer's Guide\\
       \large Bug-Bite Classification via Tropical Cyclone Transfer Learning\\
       Braude College of Engineering --- Software Engineering Capstone}
\author{Ishay Yulzary, Nathan Trostianitser}
\date{}
\maketitle
\tableofcontents
\newpage
% -------------------------------------------------------------------------

\section{Maintainer's Guide}
\label{app:maintainer-guide}

This guide is addressed to a software engineer or researcher responsible for
keeping this system operational, extending the research, or advancing the
prototype toward a deployable product. It covers the operating environment,
installation of all custom software, the training pipeline, and procedures
for updating and extending the system.

% ============================================================
\subsection{System Overview}
\label{sec:mg-overview}
% ============================================================

\subsubsection{What the System Does}

The system classifies a photograph of a bug bite into one of five categories:
ant bite, bed bug bite, mosquito bite, spider bite, or tick/flea bite. It does
this by running the photograph through an ensemble of three neural network
models and produces visual explanations showing which image regions most
influenced the prediction.

The inference interface is \texttt{Stacked\_Model\_Colab.ipynb}, a
self-contained Jupyter notebook designed for Google Colab. The full training
and evaluation pipeline runs via \texttt{miscellaneous\_code/run\_experiments.py}
on a dedicated GPU machine.


% ============================================================
\subsection{Repository Structure}
\label{sec:mg-repo}
% ============================================================

\begin{description}
    \item[\texttt{notebooks/Multiclass\_Classification.ipynb}]
          Primary experimental notebook. Covers all four pipeline stages:
          control training, TC pre-training, feature map and GradCAM
          evaluation, and the full XAI suite. Intended for interactive
          exploration or one-shot reproduction on a GPU instance.

    \item[\texttt{Stacked\_Model\_Colab.ipynb}]
          Self-contained inference and visualisation notebook for Google Colab
          free tier. Loads trained weights from Google Drive, runs ensemble
          inference on a single image, and produces side-by-side Control vs.\
          TC GradCAM and LIME outputs.

    \item[\texttt{scripts/}]
          Standalone Python scripts, one per pipeline stage, designed for
          unattended server execution via \texttt{run\_experiments.py}.
          \begin{itemize}
              \item \texttt{shared.py} --- shared imports, dataset paths,
                    metrics helpers, and GradCAM utilities.
              \item \texttt{01\_train\_control.py} --- trains the Control
                    (ImageNet $\rightarrow$ bug-bite) models.
              \item \texttt{02\_train\_tc.py} --- runs cyclone pre-training
                    then TC $\rightarrow$ bug-bite fine-tuning.
              \item \texttt{03\_evaluate.py} --- generates feature map
                    comparisons and GradCAM overlays.
              \item \texttt{04\_xai.py} --- computes ensemble metrics,
                    silhouette scores, t-SNE projections, LIME, SHAP, and
                    DiCE outputs.
          \end{itemize}

    \item[\texttt{miscellaneous\_code/pytorch\_utils.py}]
          Custom PyTorch training library imported by all scripts and the
          main notebook. See Section~\ref{sec:mg-pytorch-utils}.

    \item[\texttt{miscellaneous\_code/run\_experiments.py}]
          Automated hyperparameter sweep runner.
          See Section~\ref{sec:mg-sweep}.

    \item[\texttt{miscellaneous\_code/fetch\_cyclone\_dataset.py}]
          Downloads and labels the HURSAT-B1 cyclone dataset from NOAA.
          See Section~\ref{sec:mg-data-cyclone}.

    \item[\texttt{miscellaneous\_code/cyclone\_preprocessing.py}]
          Removes grid-line artefacts introduced by HURSAT-B1 image tiling
          before training.

    \item[\texttt{Bug-Data/Multiclass\_Bug\_data/}]
          Bug-bite dataset: \texttt{train/\{class\}/} and
          \texttt{val/\{class\}/} subdirectories, one per class.

    \item[\texttt{results/}]
          Auto-generated by the sweep runner. Contains checkpoints,
          per-run metric JSONs, XAI outputs, and plots.

    \item[\texttt{deploy/}]
          Legacy Streamlit demo (binary classifier, outdated). Not part
          of the active pipeline; retained for reference only.
\end{description}

% ============================================================
\subsection{Operating Environment}
\label{sec:mg-env}
% ============================================================

\subsubsection{Hardware Requirements}

\textbf{A dedicated GPU is required for training.} The full 63-configuration
sweep was run on the following hardware:

\begin{itemize}
    \item \textbf{CPU:} Intel Core i7-9700KF @ 5.10\,GHz all-core OC
          (8 cores / 8 threads)
    \item \textbf{RAM:} 32\,GB
    \item \textbf{GPU:} NVIDIA RTX A5000 (24\,GB GDDR6 VRAM)
\end{itemize}

The full sweep took approximately 100--130 hours of continuous GPU time on
this system. A GPU with at least 16\,GB VRAM is recommended; the default
batch size of 4 allows training on GPUs with as little as 8\,GB VRAM at the
cost of slower convergence. InceptionV3 phase-1 training may require a
reduced batch size on GPUs below 24\,GB VRAM (see
Section~\ref{sec:mg-troubleshoot-oom}).

Google Colab free tier is not suitable for the full sweep due to session time
limits. A Colab Pro instance with a high-RAM GPU runtime is the minimum viable
option for partial retraining. Inference and XAI visualisation can be run on
CPU or a smaller GPU.

\subsubsection{Software Dependencies}

\textbf{Inference (Colab):} The inference notebook installs its own
dependencies at runtime in Section 1. No manual installation is required.

\textbf{Training (local or server):}

\begin{lstlisting}[language=bash]
# Core training dependencies
pip install torch torchvision timm numpy Pillow opencv-python \
            scikit-learn matplotlib tqdm

# Optional XAI packages
pip install shap lime scikit-image dice-ml

# Cyclone dataset fetch
pip install netCDF4 requests pandas tqdm
\end{lstlisting}

There is no pinned \texttt{requirements.txt} for the training environment.
The project was developed and tested with PyTorch 2.x and timm 0.9.x.
Contact \texttt{ishay.yulzary@gmail.com} for the exact environment
specification used during the published experiments.

\subsubsection{Dataset Paths}

Dataset paths are configured at the top of
\texttt{miscellaneous\_code/run\_experiments.py} and in
\texttt{scripts/shared.py}. They default to WSL native paths:

\begin{lstlisting}[language=Python]
BUG_TRAIN = '/home/test/bug_data/train'
BUG_VAL   = '/home/test/bug_data/val'
TC_TRAIN  = '/home/test/cyclone_data_split/train'
TC_VAL    = '/home/test/cyclone_data_split/val'
\end{lstlisting}

Update these constants to match your environment before running training.

% ============================================================
\subsection{Model Weights and Artifacts}
\label{sec:mg-weights}
% ============================================================

Trained model weights are not stored in the repository. Two archives are
available from \texttt{ishay.yulzary@gmail.com}:

\begin{itemize}
    \item \texttt{project\_book\_results.tar} --- the selected best-performing
          checkpoints used in the inference demo (one Control config, one TC
          config, three architectures each).
    \item \textbf{Complete sweep archive} --- all 70 experiment configurations
          (63 TC + 7 Control seeds), covering every trained checkpoint from
          the published experiments. Request this for full evaluation
          reproduction or sweep extension.
\end{itemize}

Weight subdirectories within the demo archive:
\begin{itemize}
    \item Control: \texttt{project\_book\_results/s3\_ls0.20\_lr1.0e-05/}
    \item TC: \texttt{project\_book\_results/s6\_ls0.10\_lr5.0e-06/}
\end{itemize}

\textbf{Updating the weights archive after retraining:}
New weight files are written to \texttt{results/checkpoints/}. Package the
relevant checkpoint directories into a new
\\\texttt{project\_book\_results.tar} and upload it to Google Drive, replacing
the previous file.


% ============================================================
\subsection{Troubleshooting}
\label{sec:mg-troubleshoot}
% ============================================================

\subsubsection{GPU out-of-memory error during training}
\label{sec:mg-troubleshoot-oom}

\textbf{Symptom:} \texttt{torch.cuda.OutOfMemoryError} during
\texttt{run\_experiments.py}.

\textbf{Resolution:}
\begin{enumerate}
    \item Reduce the batch size for the affected architecture in the
          \texttt{MODELS} list in \texttt{run\_experiments.py}.
          Relevant columns: \texttt{cyc\_p1\_bs} (cyclone phase-1 batch
          size) and \texttt{bug\_bs} (bug-bite batch size).
    \item For InceptionV3, the phase-1 batch size may need to be reduced
          to 64 or below on GPUs with less than 24\,GB VRAM.
    \item Re-run the sweep. Configurations that completed before the
          crash are not repeated.
\end{enumerate}

\subsubsection{NOAA download failures during cyclone dataset fetch}

\textbf{Symptom:} Repeated warnings such as
\begin{lstlisting}
! Index fetch failed for <year>: ...
\end{lstlisting}

\textbf{Resolution:}
\begin{enumerate}
    \item Re-run the fetch script. Downloads are resumable; already-saved
          images are not re-downloaded.
    \item Reduce the worker count if the server rate-limits the connection:
\begin{lstlisting}[language=bash]
python miscellaneous_code/fetch_cyclone_dataset.py --workers 2
\end{lstlisting}
    \item If NOAA servers are unavailable, contact
          \texttt{ishay.yulzary@gmail.com} to obtain the pre-built dataset
          archive directly.
\end{enumerate}

\subsubsection{Training sweep shows no progress after resuming}

\textbf{Symptom:} Re-running \texttt{run\_experiments.py} prints
\texttt{All configs already completed} but results appear incomplete.

\textbf{Resolution:}
\begin{enumerate}
    \item Run \texttt{python miscellaneous\_code/run\_experiments.py
          --status} to view the current state of each configuration.
    \item If all 63 configurations report \texttt{completed} but results
          are missing, \\\texttt{results/experiment\_results.json} may be
          corrupted. Delete it and re-run; the sweep will re-evaluate all
          existing checkpoints without retraining.
    \item If fewer than 63 configurations are completed, pending ones
          resume automatically on the next run.
\end{enumerate}

% ============================================================
\subsection{Contacts and Resources}
\label{sec:mg-contacts}
% ============================================================

\begin{description}
    \item[Primary contact] Ishay Yulzary ---
          \texttt{ishay.yulzary@gmail.com}.
          Contact for: weight archive access, dataset access, environment
          specification, and questions not covered in this guide.

    \item[Repository] Source code, notebooks, and bug-bite dataset are
          maintained in the project GitHub repository.

    \item[Data sources] The cyclone dataset is derived from:
          \begin{itemize}
              \item HURSAT-B1:
                    \\\url{https://www.ncei.noaa.gov/products/hurricane-satellite-data}
              \item IBTrACS:
                    \\\url{https://www.ncei.noaa.gov/products/international-best-track-archive}
          \end{itemize}
          Both are freely available from NOAA.

    \item[Model library] Architectures are loaded via the \texttt{timm}
          library (\url{https://github.com/huggingface/pytorch-image-models}).
          Pre-trained ImageNet weights are downloaded automatically by
          \texttt{timm} on first use.
\end{description}

% --- End of appendix (remove if standalone) ---
\end{document}

%
%  Or compile standalone:
%    Uncomment the preamble block below, and add \end{document} at the end.
% ============================================================

% --- Standalone preamble (remove if including in main document) ----------
\documentclass[12pt,a4paper]{article}
\usepackage[utf8]{inputenc}
\usepackage[T1]{fontenc}
\usepackage{geometry}
\geometry{margin=2.5cm}
\usepackage{hyperref}
\usepackage{listings}
\usepackage{xcolor}
\usepackage{booktabs}
\usepackage{array}
\usepackage{enumitem}
\usepackage{parskip}
\lstset{basicstyle=\ttfamily\small, breaklines=true, frame=single,
        backgroundcolor=\color{gray!8}}
\begin{document}
\title{Maintainer's Guide\\
       \large Bug-Bite Classification via Tropical Cyclone Transfer Learning\\
       Braude College of Engineering --- Software Engineering Capstone}
\author{Ishay Yulzary, Nathan Trostianitser}
\date{}
\maketitle
\tableofcontents
\newpage
% -------------------------------------------------------------------------

\section{Maintainer's Guide}
\label{app:maintainer-guide}

This guide is addressed to a software engineer or researcher responsible for
keeping this system operational, extending the research, or advancing the
prototype toward a deployable product. It covers the operating environment,
installation of all custom software, the training pipeline, and procedures
for updating and extending the system.

% ============================================================
\subsection{System Overview}
\label{sec:mg-overview}
% ============================================================

\subsubsection{What the System Does}

The system classifies a photograph of a bug bite into one of five categories:
ant bite, bed bug bite, mosquito bite, spider bite, or tick/flea bite. It does
this by running the photograph through an ensemble of three neural network
models and produces visual explanations showing which image regions most
influenced the prediction.

The inference interface is \texttt{Stacked\_Model\_Colab.ipynb}, a
self-contained Jupyter notebook designed for Google Colab. The full training
and evaluation pipeline runs via \texttt{miscellaneous\_code/run\_experiments.py}
on a dedicated GPU machine.


% ============================================================
\subsection{Repository Structure}
\label{sec:mg-repo}
% ============================================================

\begin{description}
    \item[\texttt{notebooks/Multiclass\_Classification.ipynb}]
          Primary experimental notebook. Covers all four pipeline stages:
          control training, TC pre-training, feature map and GradCAM
          evaluation, and the full XAI suite. Intended for interactive
          exploration or one-shot reproduction on a GPU instance.

    \item[\texttt{Stacked\_Model\_Colab.ipynb}]
          Self-contained inference and visualisation notebook for Google Colab
          free tier. Loads trained weights from Google Drive, runs ensemble
          inference on a single image, and produces side-by-side Control vs.\
          TC GradCAM and LIME outputs.

    \item[\texttt{scripts/}]
          Standalone Python scripts, one per pipeline stage, designed for
          unattended server execution via \texttt{run\_experiments.py}.
          \begin{itemize}
              \item \texttt{shared.py} --- shared imports, dataset paths,
                    metrics helpers, and GradCAM utilities.
              \item \texttt{01\_train\_control.py} --- trains the Control
                    (ImageNet $\rightarrow$ bug-bite) models.
              \item \texttt{02\_train\_tc.py} --- runs cyclone pre-training
                    then TC $\rightarrow$ bug-bite fine-tuning.
              \item \texttt{03\_evaluate.py} --- generates feature map
                    comparisons and GradCAM overlays.
              \item \texttt{04\_xai.py} --- computes ensemble metrics,
                    silhouette scores, t-SNE projections, LIME, SHAP, and
                    DiCE outputs.
          \end{itemize}

    \item[\texttt{miscellaneous\_code/pytorch\_utils.py}]
          Custom PyTorch training library imported by all scripts and the
          main notebook. See Section~\ref{sec:mg-pytorch-utils}.

    \item[\texttt{miscellaneous\_code/run\_experiments.py}]
          Automated hyperparameter sweep runner.
          See Section~\ref{sec:mg-sweep}.

    \item[\texttt{miscellaneous\_code/fetch\_cyclone\_dataset.py}]
          Downloads and labels the HURSAT-B1 cyclone dataset from NOAA.
          See Section~\ref{sec:mg-data-cyclone}.

    \item[\texttt{miscellaneous\_code/cyclone\_preprocessing.py}]
          Removes grid-line artefacts introduced by HURSAT-B1 image tiling
          before training.

    \item[\texttt{Bug-Data/Multiclass\_Bug\_data/}]
          Bug-bite dataset: \texttt{train/\{class\}/} and
          \texttt{val/\{class\}/} subdirectories, one per class.

    \item[\texttt{results/}]
          Auto-generated by the sweep runner. Contains checkpoints,
          per-run metric JSONs, XAI outputs, and plots.

    \item[\texttt{deploy/}]
          Legacy Streamlit demo (binary classifier, outdated). Not part
          of the active pipeline; retained for reference only.
\end{description}

% ============================================================
\subsection{Operating Environment}
\label{sec:mg-env}
% ============================================================

\subsubsection{Hardware Requirements}

\textbf{A dedicated GPU is required for training.} The full 63-configuration
sweep was run on the following hardware:

\begin{itemize}
    \item \textbf{CPU:} Intel Core i7-9700KF @ 5.10\,GHz all-core OC
          (8 cores / 8 threads)
    \item \textbf{RAM:} 32\,GB
    \item \textbf{GPU:} NVIDIA RTX A5000 (24\,GB GDDR6 VRAM)
\end{itemize}

The full sweep took approximately 100--130 hours of continuous GPU time on
this system. A GPU with at least 16\,GB VRAM is recommended; the default
batch size of 4 allows training on GPUs with as little as 8\,GB VRAM at the
cost of slower convergence. InceptionV3 phase-1 training may require a
reduced batch size on GPUs below 24\,GB VRAM (see
Section~\ref{sec:mg-troubleshoot-oom}).

Google Colab free tier is not suitable for the full sweep due to session time
limits. A Colab Pro instance with a high-RAM GPU runtime is the minimum viable
option for partial retraining. Inference and XAI visualisation can be run on
CPU or a smaller GPU.

\subsubsection{Software Dependencies}

\textbf{Inference (Colab):} The inference notebook installs its own
dependencies at runtime in Section 1. No manual installation is required.

\textbf{Training (local or server):}

\begin{lstlisting}[language=bash]
# Core training dependencies
pip install torch torchvision timm numpy Pillow opencv-python \
            scikit-learn matplotlib tqdm

# Optional XAI packages
pip install shap lime scikit-image dice-ml

# Cyclone dataset fetch
pip install netCDF4 requests pandas tqdm
\end{lstlisting}

There is no pinned \texttt{requirements.txt} for the training environment.
The project was developed and tested with PyTorch 2.x and timm 0.9.x.
Contact \texttt{ishay.yulzary@gmail.com} for the exact environment
specification used during the published experiments.

\subsubsection{Dataset Paths}

Dataset paths are configured at the top of
\texttt{miscellaneous\_code/run\_experiments.py} and in
\texttt{scripts/shared.py}. They default to WSL native paths:

\begin{lstlisting}[language=Python]
BUG_TRAIN = '/home/test/bug_data/train'
BUG_VAL   = '/home/test/bug_data/val'
TC_TRAIN  = '/home/test/cyclone_data_split/train'
TC_VAL    = '/home/test/cyclone_data_split/val'
\end{lstlisting}

Update these constants to match your environment before running training.

% ============================================================
\subsection{Model Weights and Artifacts}
\label{sec:mg-weights}
% ============================================================

Trained model weights are not stored in the repository. Two archives are
available from \texttt{ishay.yulzary@gmail.com}:

\begin{itemize}
    \item \texttt{project\_book\_results.tar} --- the selected best-performing
          checkpoints used in the inference demo (one Control config, one TC
          config, three architectures each).
    \item \textbf{Complete sweep archive} --- all 70 experiment configurations
          (63 TC + 7 Control seeds), covering every trained checkpoint from
          the published experiments. Request this for full evaluation
          reproduction or sweep extension.
\end{itemize}

Weight subdirectories within the demo archive:
\begin{itemize}
    \item Control: \texttt{project\_book\_results/s3\_ls0.20\_lr1.0e-05/}
    \item TC: \texttt{project\_book\_results/s6\_ls0.10\_lr5.0e-06/}
\end{itemize}

\textbf{Updating the weights archive after retraining:}
New weight files are written to \texttt{results/checkpoints/}. Package the
relevant checkpoint directories into a new
\\\texttt{project\_book\_results.tar} and upload it to Google Drive, replacing
the previous file.


% ============================================================
\subsection{Troubleshooting}
\label{sec:mg-troubleshoot}
% ============================================================

\subsubsection{GPU out-of-memory error during training}
\label{sec:mg-troubleshoot-oom}

\textbf{Symptom:} \texttt{torch.cuda.OutOfMemoryError} during
\texttt{run\_experiments.py}.

\textbf{Resolution:}
\begin{enumerate}
    \item Reduce the batch size for the affected architecture in the
          \texttt{MODELS} list in \texttt{run\_experiments.py}.
          Relevant columns: \texttt{cyc\_p1\_bs} (cyclone phase-1 batch
          size) and \texttt{bug\_bs} (bug-bite batch size).
    \item For InceptionV3, the phase-1 batch size may need to be reduced
          to 64 or below on GPUs with less than 24\,GB VRAM.
    \item Re-run the sweep. Configurations that completed before the
          crash are not repeated.
\end{enumerate}

\subsubsection{NOAA download failures during cyclone dataset fetch}

\textbf{Symptom:} Repeated warnings such as
\begin{lstlisting}
! Index fetch failed for <year>: ...
\end{lstlisting}

\textbf{Resolution:}
\begin{enumerate}
    \item Re-run the fetch script. Downloads are resumable; already-saved
          images are not re-downloaded.
    \item Reduce the worker count if the server rate-limits the connection:
\begin{lstlisting}[language=bash]
python miscellaneous_code/fetch_cyclone_dataset.py --workers 2
\end{lstlisting}
    \item If NOAA servers are unavailable, contact
          \texttt{ishay.yulzary@gmail.com} to obtain the pre-built dataset
          archive directly.
\end{enumerate}

\subsubsection{Training sweep shows no progress after resuming}

\textbf{Symptom:} Re-running \texttt{run\_experiments.py} prints
\texttt{All configs already completed} but results appear incomplete.

\textbf{Resolution:}
\begin{enumerate}
    \item Run \texttt{python miscellaneous\_code/run\_experiments.py
          --status} to view the current state of each configuration.
    \item If all 63 configurations report \texttt{completed} but results
          are missing, \\\texttt{results/experiment\_results.json} may be
          corrupted. Delete it and re-run; the sweep will re-evaluate all
          existing checkpoints without retraining.
    \item If fewer than 63 configurations are completed, pending ones
          resume automatically on the next run.
\end{enumerate}

% ============================================================
\subsection{Contacts and Resources}
\label{sec:mg-contacts}
% ============================================================

\begin{description}
    \item[Primary contact] Ishay Yulzary ---
          \texttt{ishay.yulzary@gmail.com}.
          Contact for: weight archive access, dataset access, environment
          specification, and questions not covered in this guide.

    \item[Repository] Source code, notebooks, and bug-bite dataset are
          maintained in the project GitHub repository.

    \item[Data sources] The cyclone dataset is derived from:
          \begin{itemize}
              \item HURSAT-B1:
                    \\\url{https://www.ncei.noaa.gov/products/hurricane-satellite-data}
              \item IBTrACS:
                    \\\url{https://www.ncei.noaa.gov/products/international-best-track-archive}
          \end{itemize}
          Both are freely available from NOAA.

    \item[Model library] Architectures are loaded via the \texttt{timm}
          library (\url{https://github.com/huggingface/pytorch-image-models}).
          Pre-trained ImageNet weights are downloaded automatically by
          \texttt{timm} on first use.
\end{description}

% --- End of appendix (remove if standalone) ---
\end{document}
